\newpage
\section{Ponto de projeto}

Adotou-se o mesmo ponto de projeto do Lab 03, a condição de cruzeiro com o
peso médio da fase, $W = 229\,669{,}3$ kgf, que corresponde a cerca de 44\%
da capacidade de combustível com 100\% de carga paga. A escolha mantém o
$C_L$ de projeto igual ao usado na otimização dos perfis, de modo que as
análises deste trabalho avaliam a asa exatamente na condição para a qual as
seções foram desenhadas. A Tabela~\ref{tab:ponto_projeto} reúne os dados
pedidos no roteiro.

\begin{table}[H]
\centering
\caption{Dados do ponto de projeto.}\label{tab:ponto_projeto}
\renewcommand{\arraystretch}{1.25}
\begin{tabular}{lcl}
\toprule
Parâmetro & Valor & Explicação \\
\midrule
$W_0$ & 2.857.571 & Peso máximo de decolagem da aeronave [N] \\
$W$ & 2.253.056 & Peso da aeronave no ponto de projeto [N] \\
$h$ & 10.668 & Altitude no ponto de projeto [m] \\
$\rho_\infty$ & 0,38046 & Densidade do ar no ponto de projeto [kg/m$^3$] \\
$a_\infty$ & 296,587 & Velocidade do som no ponto de projeto [m/s] \\
$M$ & 0,85 & Mach no ponto de projeto \\
$V$ & 252,1 & Velocidade no ponto de projeto [m/s] \\
$C_L$ & 0,5053 & $C_L$ no ponto de projeto \\
$S_{\mathrm{ref}}$ & 368,833 & Área de referência [m$^2$] \\
\bottomrule
\end{tabular}
\end{table}

O coeficiente de sustentação segue diretamente do equilíbrio vertical na
condição de cruzeiro:
\begin{equation*}
    C_L = \frac{W}{\tfrac{1}{2}\,\rho_\infty\,V^2\,S_{\mathrm{ref}}}
        = 0{,}5053.
\end{equation*}
Este é o valor usado como meta de sustentação (\texttt{a c 0.5053}) em todas
as simulações de cruzeiro do AVL neste relatório.
